\section{What the Schema Says About Your Types}
\label{sec:ch02-what-the-schema-says-about-your-types}

\index{schema!export}%
Nothing so far required you to look at the schema, and that is exactly the
problem. Ask the program for it. \code{RunWithGraphQLCommands} in
\code{Program.cs} is what makes the subcommand available, and the double dash
separates your arguments from the ones \code{dotnet run} wants.

\begin{minted}{text}
dotnet run -- schema export --output schema.graphql
\end{minted}

It writes two files. The second, \code{schema-settings.json}, is a workspace
file for the IDE in section~\ref{sec:ch02-the-development-loop} and records
whatever address the export run happened to bind, which is not the one your
launch profile pins. Nothing in this book uses it. The first is the schema, and
these are its type definitions:

\begin{graphqlsdl}
schema {
  query: Query
}

type Query {
  "Every session on the schedule, earliest first."
  sessions: [Session!]! @cost(weight: "10")
  "One session by its identifier, or null if no session has it."
  sessionById(id: Int!): Session @cost(weight: "10")
}

"A talk on the conference program."
type Session {
  "The person giving this session."
  speaker: Speaker
  "Stable identifier for this session."
  id: Int!
  "The line that goes in the program."
  title: String!
  "The longer text, if there is one."
  abstract: String
  "When it starts, in conference time."
  startsAt: DateTime!
  "How long the slot is."
  durationMinutes: Int!
}

"A person giving one or more sessions."
type Speaker {
  "Stable identifier for this speaker."
  id: Int!
  "The name to print in the program."
  name: String!
  "A sentence or two, if the speaker sent any."
  bio: String
}
\end{graphqlsdl}

The rest of the file is the two things those definitions lean on, neither of
which you wrote:

\begin{graphqlsdl}[fontsize=\footnotesize, breakanywhere]
"The `DateTime` scalar type represents a date and time with time zone offset information."
scalar DateTime
  @specifiedBy(url: "https://scalars.graphql.org/chillicream/date-time.html")

"The purpose of the `cost` directive is to define a `weight` for GraphQL types, fields, and arguments. Static analysis can use these weights when calculating the overall cost of a query or response."
directive @cost(
  "The `weight` argument defines what value to add to the overall cost for every appearance, or possible appearance, of a type, field, argument, etc."
  weight: String!
) on
  | SCALAR
  | OBJECT
  | FIELD_DEFINITION
  | ARGUMENT_DEFINITION
  | ENUM
  | INPUT_FIELD_DEFINITION
\end{graphqlsdl}

Read that against the C\# and every line of it is a rule you can now
predict.

\index{naming!Get prefix}%
\index{nullability!mapped from C\#}%
\index{scalar!DateTime}%
\begin{center}
\begin{tabular}{@{}ll@{}}
\toprule
\textrm{What you wrote} & \textrm{What the schema says} \\
\midrule
\code{GetSessions()}            & \code{sessions}              \\
\code{GetSessionById(int id)}   & \code{sessionById(id: Int!)} \\
\code{string Title}             & \code{title: String!}        \\
\code{string? Abstract}         & \code{abstract: String}      \\
\code{DateTimeOffset StartsAt}  & \code{startsAt: DateTime!}   \\
\bottomrule
\end{tabular}
\end{center}

The \code{Get} prefix is stripped and the first letter is lowercased, so
the C\# convention and the GraphQL convention are reconciled without either
side giving way. An argument keeps its name and its type.

The second pair is the most consequential mapping in Hot Chocolate and the
one people discover last. Your nullable reference type annotations are the
schema's non-null annotations. A client's code generator reads
\code{String!} and emits a non-optional field, so a question mark you did
not type in C\# becomes a crash in somebody else's application. It is also
why \code{sessionById} returns \code{Session} rather than \code{Session!},
which is the whole reason that asking for a session that is not there is an
answer instead of an error. Chapter~\ref{ch:null-blast-radius} is about a
field marked non-null resolving to null anyway, across a seam, and it
collects on every decision made here.

The last pair introduces a scalar the server defines itself. \code{DateTime} is
not a built-in, which is why the export carries a definition for it and an
\code{@specifiedBy} pointing at a written specification: a client that has
never met the type has to be told what is in the string.

\index{schema!descriptions from XML comments}%
Every field carries a description, and none of them was written in a schema
file. A \code{<summary>} on a method, and a \code{<param>} on a positional
record's own doc comment, both become the description of the field they
name. The second is worth knowing, because there is no reason to expect
that a tag documenting a constructor parameter would end up describing a
schema field. Delete \code{GenerateDocumentationFile} from the project file
and every one of these quoted lines disappears, silently, leaving the
schema undocumented and the build clean.

\index{cost directive}%
Then there is \code{@cost(weight: "10")}, which appears on both fields and
which nobody asked for. Hot Chocolate applies a cost weight to fields with no
configuration at all, so complexity analysis is present in the schema from the
first build. It is not a version~16 novelty either; the same directive comes
out of a version~14 build of these same files. Chapter~\ref{ch:operations} is
where the numbers on it start mattering.

\index{schema!field order}%
One detail is a curiosity rather than a rule, and it is worth naming so you do
not go looking for meaning in it. \code{speaker} is listed first on
\code{Session}, ahead of \code{id} and \code{title}, because the type extension
is applied before the record's own members are read. Field order in a schema
carries no semantics. It does mean that a schema file checked into version
control will reorder itself if you move a resolver, which matters the day
schema diffs become a gate rather than a curiosity, in
chapter~\ref{ch:breaking-changes}.

\index{schema!leaked fields}%
Now the field that is not there. \code{SpeakerId} is a public property on a
public record, and without \code{[GraphQLIgnore]} it would be
\code{speakerId: Int!} in the schema, between \code{abstract} and
\code{startsAt}, published to every client, forever. Nothing warned me. Nothing
would have warned you. This is the cost of the sentence this chapter opened
with: the schema is derived from your types, so a member you added for the
compiler is an API commitment until you say otherwise, and the direction of the
default is publish.

I would rather have the opposite default, and I do not get it. What I do
instead is read the exported schema after every change that adds a member, and
treat \code{schema.graphql} as a reviewed file rather than a build artifact.
